\documentclass[10pt,a4paper]{amsart}
\usepackage[margin=19mm]{geometry}
\usepackage{amsmath,amssymb,mathtools}
\usepackage[scr=rsfs,cal=euler]{mathalfa}
\usepackage{xfrac}
\usepackage[unicode,hidelinks,bookmarks=false]{hyperref}
\newcommand{\ie}{i.e.\ }
\newcommand{\cf}{cf.\ }
\newcommand{\resp}{respectively\ }
\newcommand{\Subsection}{\subsection}
\providecommand{\nobreakdash}{\nobreak\hbox{-}}
\setlength{\emergencystretch}{2em}
\multlinegap0pt
\usepackage{bm}
\usepackage{bbm}
\usepackage{dsfont}
\usepackage{xspace}
\usepackage{cancel}
\usepackage{mathtools}
\usepackage{cleveref}
\usepackage{amsthm}
\makeatletter
\@ifundefined{thm}{\newtheorem{thm}{Thm}}{}
\@ifundefined{proposition}{\newtheorem{proposition}[thm]{Proposition}}{}
\@ifundefined{theorem}{\newtheorem{theorem}[thm]{Theorem}}{}
\makeatother
\newcommand{\Q}{\mathbb{Q}}
\newcommand{\V}{\mathcal{V}}
\newcommand{\rk}{\operatorname{rk}}
\title[The asymptotic size of finite irreducible semigroups of rati]{The asymptotic size of finite irreducible semigroups of rational matrices}
\author{Stefan Kiefer, Andrew Ryzhikov}
\date{Source: arXiv:2601.01236v2 (CC BY 4.0)}
\begin{document}
\maketitle

\noindent\emph{Source and licence.} The asymptotic size of finite irreducible semigroups of rational matrices. Version arXiv:2601.01236v2 (CC BY 4.0), distributed under the Creative Commons Attribution 4.0 International licence, as recorded in the arXiv licence metadata for this exact version and shown on the abstract page https://arxiv.org/abs/2601.01236v2. Full rights notice: licenses/arxiv-2601.01236-CC-BY-4.0.txt.\par\smallskip

\noindent\textit{Editorial scope.} Counted unit: the Proposition whose source label is \texttt{prop:min-rank-diameter} in the article (source0.tex lines 1326--1329, source label \texttt{prop:min-rank-diameter}), delivered together with the article's own complete original proof of it (source0.tex lines 1330--1394, one \texttt{proof} environment of 65 lines). No mathematics is written, repaired or re-derived: statement and proof are the article's own text line for line. Required context retained uncounted: thm:main-bound (lines 1275--1277); prop:diameter-size-bound (lines 1308--1312). Every definition, display and label of those lines is retained unchanged. Omitted from this package and not redistributed: 1--1274, 1278--1307, 1313--1325, 1395--1958. Nothing inside a retained excerpt is omitted or altered: every retained body is delivered exactly as the article's lines are written, so no omission receipt of any kind accompanies this package. Citations: the retained excerpts carry no citation command at all, so no bibliography entry of the article is transcribed and no reference is redistributed. Reference adaptation: none was needed and none was made; every reference inside the delivered unit resolves inside it, so no unresolved reference or citation is produced. Printed slips kept literal: no printed word, symbol, hypothesis or number of the counted statement or of its proof was altered. Layout and class: the article's own class is \texttt{lmcs} with packages including babel, graphicx, mathtools, amsmath, amssymb, amsfonts, caption, subcaption, xspace, xcolor, cleveref; the wrapper is the lane's pinned \texttt{amsart} editorial wrapper at 10pt on A4, which restates the theorem-like environments the retained text needs and adds the article's own macro definitions that the retained text needs (\texttt{\textbackslash Q}, \texttt{\textbackslash V}, \texttt{\textbackslash rk}), with extra wrapper packages (bm, bbm, dsfont, xspace, cancel, mathtools, cleveref). The delivered numbering is the unit's own. Assets: no retained span contains a figure environment, a graphics-inclusion command, a PDF-page inclusion, a table or a TikZ picture; no image file is copied into this package, the package declares no asset dependency and no image file is redistributed with it. Licence: version arXiv:2601.01236v2 of this article is distributed under the Creative Commons Attribution 4.0 International licence, as recorded in the arXiv licence metadata for this exact version and shown on the abstract page https://arxiv.org/abs/2601.01236v2. A full rights notice is delivered at licenses/arxiv-2601.01236-CC-BY-4.0.txt. Characters: every retained excerpt is pure ASCII, so the delivered engine, XeLaTeX with its default Latin Modern text font, reports no missing character.
\begin{theorem}\label{thm:main-bound}
Let $S \subseteq \Q^{n \times n}$ be a finite irreducible semigroup. Then $|S| \le 3^{n^2}$.
\end{theorem}

\begin{proposition} \label{prop:diameter-size-bound}
Let $S \subseteq \Q^{n \times n}$ be a finite, not necessarily irreducible, semigroup, generated by $S_0 \subseteq S$.
For any $X \in S$, its depth is at most~$|S|$.
Hence, the minimum-rank diameter is at most~$|S|$.
\end{proposition}

\begin{proposition} \label{prop:min-rank-diameter}
Let $S \subseteq \Q^{n \times n}$ be a finite, not necessarily irreducible, semigroup, generated by $S_0 \subseteq S$.
Then its minimum-rank diameter is at most~$3^{n^2}$.
\end{proposition}

\begin{proof}
We prove the proposition by (strong) induction on~$n \ge 1$.
Let $n \ge 1$ and suppose the proposition holds for all $1 \le m < n$.
Let the finite semigroup $S \subseteq \Q^{n \times n}$ be generated by~$S_0 \subseteq S$.
Let $M : \Sigma^+ \to S$ be a semigroup homomorphism with $M(\Sigma) = S_0$ and $M(\Sigma^+) = S$.

Suppose $S$~is irreducible.
By \cref{thm:main-bound} we have $|S| \le 3^{n^2}$.
Hence, by \cref{prop:diameter-size-bound} the minimum-rank diameter is at most~$3^{n^2}$.

So we can assume that $S$~is not irreducible.
Choose~$\V_1$ to be a minimal nonzero $S$-invariant subspace of~$\Q^n$; then $\{\vec{0}\} \ne \V_1 \ne \Q^n$.
Let $\V_2 \subseteq \Q^n$ be a complement of~$\V_1$ so that $\Q^n = \V_1 \oplus \V_2$.
Let $n_i \coloneqq \dim \V_i$ for $i=1,2$.
We have $n = n_1 + n_2$ with $n_1,n_2 \ge 1$.
Since $\V_1$~is $S$-invariant, in a basis adapted to $\Q^n = \V_1 \oplus \V_2$ we have
\[
M(w) \ = \ \begin{pmatrix} M_1(w) & N(w) \\ 0 & M_2(w) \end{pmatrix} \qquad \text{for all $w \in \Sigma^+$,}
\]
where $M_i(w) \in \Q^{n_i \times n_i}$ for $i = 1,2$, and $N(w) \in \Q^{n_1 \times n_2}$.

Define $S_i \coloneqq M_i(\Sigma^+)$ for $i=1,2$.
Then for $i=1,2$ the set $S_i \subseteq \Q^{n_i \times n_i}$ is a finite semigroup generated by~$M_i(\Sigma)$.
The semigroup $S_1$ is irreducible; indeed, identifying $\Q^{n_1}$ with~$\V_1$ via the chosen basis, any proper nonzero $S_1$-invariant subspace of~$\Q^{n_1}$ would correspond to a proper nonzero $S$-invariant subspace of~$\V_1$, contradicting the minimality of~$\V_1$.

It follows from the block-triangular shape above that $\rk M(w) \ge \rk M_1(w) + \rk M_2(w)$ for all $w \in \Sigma^+$.
In particular, defining $r \coloneqq \min\{\rk X \mid X \in S\}$ and $r_i \coloneqq \min\{\rk X \mid X \in S_i\}$ for $i=1,2$, we have $r \ge r_1 + r_2$.
For $i=1,2$ let $w_i \in \Sigma^+$ be shortest words with $\rk M_i(w_i) = r_i$.

Since $S_1$~is irreducible, by \cref{thm:main-bound} we have $|S_1| \le 3^{n_1^2}$, hence by \cref{prop:diameter-size-bound}, $|w_1| \le 3^{n_1^2}$.
Applying the induction hypothesis to~$S_2$ we obtain $|w_2| \le 3^{n_2^2}$.
Further,
\begin{align*}
M(w_1 w_2) \ &= \ M(w_1) M(w_2) \ = \ 
\begin{pmatrix}
    M_1(w_1) & N(w_1) \\ 0 & M_2(w_1)
\end{pmatrix}
\begin{pmatrix}
    M_1(w_2) & N(w_2) \\ 0 & M_2(w_2)
\end{pmatrix} \\
&= \ 
\underbrace{\begin{pmatrix}
    M_1(w_1) \\ 0 
\end{pmatrix}}_{\text{rank } r_1}
\begin{pmatrix}
    M_1(w_2) & N(w_2)
\end{pmatrix}
+
\begin{pmatrix}
    N(w_1) \\ M_2(w_1)
\end{pmatrix}
\underbrace{
\begin{pmatrix}
   0 & M_2(w_2)
\end{pmatrix}
}_{\text{rank } r_2}\,.
\end{align*}
Thus, using $\rk(X Y) \le \min \{\rk X, \rk Y\}$, the first summand has rank $\le r_1$, and the second summand has rank $\le r_2$.
Using $\rk(X + Y) \le \rk X + \rk Y$, we conclude that $\rk M(w_1 w_2) \le r_1 + r_2$.
Since $r_1 + r_2 \le r$, the matrix $M(w_1 w_2)$ is a minimum-rank matrix in~$S$.
Hence, the minimum-rank diameter is at most
\[
  |w_1 w_2| \ = \ |w_1| + |w_2| \ \le \ 3^{n_1^2} + 3^{n_2^2} \ \le \ 2 \cdot 3^{\max\{n_1^2, n_2^2\}} \ \le \ 3^{n_1^2 + n_2^2} \ \le \ 3^{(n_1+n_2)^2} \ = \ 3^{n^2}\,. \qedhere
\]
\end{proof}

\end{document}
